\section{Chapitre~\ref{sec:motorlearning}. L'apprentissage moteur}

La règle des moindres carrés et l'avantage d'un fil d'erreur séparé sont des théorèmes; la structure du circuit à fil d'erreur, l'affaiblissement par coïncidence, l'ajustement de la trace au retard, l'effet consécutif et le retour spontané de l'habileté sont mesurés; que le circuit entier fonctionne comme un apprentissage supervisé et construise un modèle interne est l'hypothèse dominante; les chiffres du modèle à deux processus sont calculés avec le modèle du livre.

{\footnotesize
\setlength{\tabcolsep}{3pt}\renewcommand{\arraystretch}{1.15}
\begin{longtable}{>{\raggedright}p{0.5cm}>{\raggedright}p{4.0cm}>{\raggedright}p{5.6cm}>{\raggedright}p{2.3cm}>{\raggedright\arraybackslash}p{1.7cm}}
\toprule
N\textsuperscript{o} & Affirmation & Expérience et ce qui est mesuré & Source & Statut \\
\midrule\endfirsthead
\toprule
N\textsuperscript{o} & Affirmation & Expérience et ce qui est mesuré & Source & Statut \\
\midrule\endhead
1 & La règle~\eqref{eq:lms} suit en moyenne le gradient du carré de l'erreur & Résultat de la théorie des filtres adaptatifs: une correction proportionnelle à l'entrée et à l'erreur de sortie est une descente stochastique du gradient de l'erreur quadratique moyenne. & \cite{WidrowHoff1960} & théorie \\
2 & Avec un fil d'erreur par sortie, la vitesse ne dépend pas du nombre de sorties; avec un seul nombre commun, elle baisse avec le nombre de neurones bruités (proposition~\ref{prop:teachers}) & Courbes d'apprentissage d'un réseau linéaire: l'apprentissage par perturbations aléatoires des neurones est plus lent que le gradient exact d'un facteur égal au nombre de neurones perturbés. & \cite{Werfel2005} & théorie \\
3 & Le circuit à fil d'erreur fonctionne comme un apprentissage supervisé (proposition~\ref{prop:errorwire}) & Théories déduites de la structure du circuit. Les expériences des lignes 7--10 concordent avec elles; que le circuit entier fonctionne exactement ainsi n'est pas démontré. & \cite{Marr1969, Albus1971} & hypothèse \\
4 & Couche d'expansion: environ $7\cdot10^{10}$ petits neurones \nt{GLUT}, plus que dans tout le reste du cerveau & Comptage des noyaux dans une suspension de tissu dissous: $69\cdot10^9$ neurones dans le cervelet humain sur $86\cdot10^9$ dans tout le cerveau. Stéréologie: $101\cdot10^9$ petits neurones dans le cervelet d'hommes âgés. & \cite{Azevedo2009, Andersen1992cereb} & mesuré \\
5 & Augmenter la dimension de l'entrée rend linéaire la dépendance voulue & Théorème sur le nombre de dichotomies: une somme pondérée de $n$ entrées avec seuil réalise un étiquetage arbitraire d'environ $2n$ vecteurs en position générale. Que le cervelet exploite précisément cela fait partie de l'hypothèse de la ligne~3. & \cite{Cover1965} & théorie \\
6 & Environ $3\cdot10^7$ neurones \nt{GABA} de sortie, chacun avec de l'ordre de $10^5$ entrées & Stéréologie du cervelet humain: $30.5\cdot10^6$ neurones de sortie. Microscopie électronique, rat: environ $1.75\cdot10^5$ entrées de la couche d'expansion par neurone de sortie. Le neurotransmetteur inhibiteur des neurones de sortie: dans la revue. & \cite{Andersen1992cereb, NapperHarvey1988, Ito2001} & mesuré \\
7 & Exactement un fil d'erreur par neurone de sortie, environ une fois par seconde, avec chaque fois une décharge puissante & Revue d'enregistrements: chaque neurone de sortie reçoit une seule entrée \nt{GLUT} grimpante, qui provoque une impulsion complexe à une fréquence de l'ordre de $1$~Hz. & \cite{Ito2001} & mesuré \\
8 & La probabilité de la décharge dépend de la direction de l'erreur de mouvement & Singe, région oculomotrice du cervelet, adaptation des saccades: la direction de l'erreur visuelle fixe la probabilité de l'impulsion complexe, sa grandeur fixe son instant; la décharge modifie les impulsions simples à l'essai suivant et déplace le mouvement le long de l'erreur préférée de la cellule. & \cite{Herzfeld2018} & mesuré \\
9 & La coïncidence de la décharge d'erreur avec une entrée active affaiblit cette entrée ($-\eta e_i x_j$) & Revue d'expériences: la stimulation conjointe d'une entrée de la couche d'expansion et du fil d'erreur provoque un affaiblissement durable de cette entrée; la stimulation de l'entrée seule, non. & \cite{Ito2001} & mesuré \\
10 & La durée de la trace est ajustée au retard de l'erreur: pour un retard de $100\ms$, c'est l'entrée active $100\ms$ avant qui s'affaiblit le plus & Tranches et souris vivantes: dans la région chargée de l'apprentissage oculomoteur, l'affaiblissement est étroitement accordé sur un intervalle d'environ $120\ms$ entre l'entrée et l'erreur, soit le temps que met le signal d'erreur dans cette tâche; dans une autre région, différentes cellules sont accordées sur différents intervalles. & \cite{Suvrathan2016} & mesuré \\
11 & Le circuit est un filtre adaptatif~\eqref{eq:adaptivefilter} qui apprend un modèle interne du corps & Modèle déduit de la structure du circuit. Aucune expérience directe n'a mesuré le filtre appris comme fonction de transfert du corps. & \cite{Fujita1982} & hypothèse \\
12 & La prédiction soustrait les conséquences de nos propres mouvements, c'est pourquoi on ne peut pas se chatouiller soi-même & Sujets humains: un toucher par soi-même semble moins chatouilleux que le même toucher par autrui; en IRM, le cortex somatosensoriel répond plus faiblement au toucher par soi-même, et le cervelet est moins actif quand le mouvement touche la peau. L'explication par soustraction de la prédiction est une hypothèse. & \cite{Blakemore1998} & hypothèse \\
13 & Sous une force extérieure, le raté diminue, et après sa suppression la main rate de l'autre côté & Des sujets tendent le bras en tenant la poignée d'un robot dans un champ de forces: les trajectoires redeviennent droites; après la suppression du champ, l'effet consécutif est presque l'image en miroir des premières déformations; il se transfère aussi à des régions non entraînées de l'espace. & \cite{ShadmehrMussaIvaldi1994} & mesuré \\
14 & Deux processus apprennent~\eqref{eq:tworate}; après une perturbation inverse, l'habileté revient d'elle-même & Sujets humains, champ de forces, puis bref champ inverse et essais à erreur bloquée: la correction revient d'elle-même vers l'ancienne; le modèle à deux processus le reproduit, le modèle à un processus non. & \cite{Smith2006} & mesuré \\
15 & Chiffres de la fig.~\ref{fig:tworate}: $0.84$ (lent $0.63$) en $200$ mouvements; $6$ inverses; rapide $-0.53$, lent $0.46$; retour jusqu'à $0.37$ en $40$ mouvements & Calcul avec \texttt{models/\allowbreak motor\_\allowbreak learning.py}, $A_f=0.92$, $B_f=0.1$, $A_s=0.996$, $B_s=0.02$: $0.840$ ($0.634$ et $0.205$), $6$ mouvements, $-0.531$ et $0.457$, pic de $0.371$ au $40$\textsuperscript{e} mouvement. & \texttt{models/\allowbreak motor\_\allowbreak learning.py} & modèle \\
16 & Deux processus sont nécessaires parce que le corps change à des vitesses différentes & Théorie: l'estimation optimale sous des perturbations à plusieurs durées de vie donne plusieurs filtres de constantes de temps différentes et explique le retour spontané et le réapprentissage accéléré. & \cite{Kording2007} & hypothèse \\
\bottomrule
\end{longtable}}
